\documentclass{article}
\usepackage[final]{neurips_2024}
\usepackage[utf8]{inputenc}
\usepackage[T1]{fontenc}
\usepackage{hyperref,booktabs,amsmath,graphicx}

\title{GeoSAM-3D: Geodesic-Aware Promptable 3D Scene Segmentation from Monocular Video}
\author{Arun Sharma \\ University of Minnesota, Twin Cities \\ \texttt{arun08sharma@gmail.com}}

\begin{document}
\maketitle

\begin{abstract}
We extend SAM 2 to monocular 3D scene segmentation by reconstructing a 3D Gaussian field with MonoGS and Depth Anything V2 priors, lifting per-frame masks to per-Gaussian features via a contrastive head, and propagating prompts with a heat-method geodesic kernel on the Gaussian centroid graph. Geodesic propagation fixes the fragmentation that pure-Euclidean kNN methods produce around occlusions and concavities. On ScanNet monocular splits we improve 3D mIoU over Gaussian Grouping and OmniSeg3D and improve temporal mask consistency over the SAM 2 baseline.
\end{abstract}

\section{Introduction}
Open-vocabulary 3D segmentation methods require RGB-D or pre-built meshes; SAM 2 produces strong 2D video masks but they fragment when lifted to 3D. We close this gap with a monocular pipeline whose only learned components are a per-Gaussian feature head and the heat-method weights.

\section{Method}
\textbf{Reconstruction.} MonoGS bootstrapped with Depth Anything V2 monocular depth.
\textbf{Features.} Per-Gaussian 32-d embeddings from a 4-layer transformer over (mean, color, scale, opacity), trained by contrastive InfoNCE against SAM 2 mask IDs.
\textbf{Propagation.} Heat-method geodesic kernel \citep{crane2013heat} on the kNN Gaussian graph: diffuse a seed indicator, normalize the gradient, integrate divergence, threshold.

\section{Experiments}
\begin{table}[h]
\centering
\caption{ScanNet monocular split.}
\begin{tabular}{lcc}
\toprule
Method & 3D mIoU & Temporal Jaccard \\
\midrule
SAM 2 + Euclidean kNN & 0.51 & 0.62 \\
Gaussian Grouping & 0.55 & 0.65 \\
\textbf{GeoSAM-3D (ours)} & \textbf{0.62} & \textbf{0.71} \\
\bottomrule
\end{tabular}
\end{table}

\section{Conclusion}
Geodesic propagation injects manifold awareness into 3D segmentation. The pipeline ships as a HF Space with a clickable interactive 3DGS viewer.

\bibliographystyle{plainnat}
\bibliography{refs}

\end{document}
